% ============================================================================
%  Section VIII — Conclusion  (Measurement, Elsevier)
% ============================================================================

\section{Conclusion}
\label{sec:conclusion}

This paper asked how the diagnostic performance of a low-cost, low-rate
smartphone-MEMS sensing chain should be measured, and what that measurement then
says about deploying such a chain on an induction machine. Using a public
$100$~Hz dataset, the near-perfect accuracies reported under random-split
evaluation were shown to be an artefact of the evaluation protocol: under a
recording-wise protocol the accuracy of standard classifiers falls by
approximately $41$ to $47$ percentage points, with comparable reductions under
cross-speed and cross-load shifts.

The more useful result is the attribution of that gap. Removing window overlap
changes nothing, and removing temporal proximity between windows changes nothing
either; the entire $46.7$-point collapse occurs only when whole recordings are
held out. Sharing samples across the partition boundary is therefore not the
operative mechanism, and the precaution most commonly taken against leakage in
this literature---non-overlapping windows---would have provided no protection
whatever. Where each operating condition is captured in one continuous
acquisition, the recording is the unit that must be separated.

A second finding concerns what a sampling-rate sweep measures. Decimation to
$25$~Hz improves leakage-free accuracy even though it removes every shaft
rotational frequency present in the dataset. Band-limiting to the same range at
the original rate is in fact mildly harmful, the improvement follows from the
smaller number of samples per analysis window, and omitting the anti-alias filter
costs $24.9$ points. The time-domain feature set is not invariant to the sampling
rate, so a rate sweep conducted with a fixed feature set measures the observable
band and the sample-count sensitivity of the features at once. The practical
guidance survives---a resource-frugal configuration of $25$~Hz, three raw channels
and a $4$~s window reaches $75.0\pm11.4\%$ against $53.9\pm12.5\%$ for the
full-rate baseline, a paired improvement of $21.1$ points ($p=0.020$), while
cutting data volume roughly eight-fold---but the explanation for it changes, and
the intermediate step of rate reduction alone is not significant at the recording
level. Repeating the selection inside a nested recording-wise loop, so that the
recordings used for estimation are never seen by the selector, returns the same
configuration in $13$ of $18$ folds but lowers its estimated performance to
$71.0\%$, a selection bias of $6.0$ points that is smaller than the effect claimed
yet comparable to the spacing between neighbouring configurations.

Reporting practice matters as much as the protocol. Scores were computed per
held-out recording and summarised with intervals clustered on that unit, giving
intervals of roughly $\pm12$ points; the same data summarised over repeated
re-partitionings would have produced intervals about three times narrower and
significance far stronger than the evidence supports. For a single-machine study
in which each fault class is one physical specimen, that is the honest width.

Future work should extend the evaluation across machines of different ratings and
manufacturers, and across repeated installations of the same fault, since
mounting is known to be a first-order influence quantity and is the confound this
design cannot address. Complementing the analysis with direct measurement of
on-node inference latency and energy would replace the calculated deployment
quantities with measured ones. The evaluation protocol set out here is intended
to serve as a deployment-honest baseline for those efforts.
